\section{The constructible inner model}
\label{sec:annex-p06-constructible-inner-model-completion}

\begin{annexitem}{annex:p06-collapse-complete}{thm:p06-collapse-complete}{The collapse construction}
\begin{proof}
For \(a\in X\), let \(D(a)\) consist of \(a\) and the points reached
by finite chains of \(R\)-predecessors.
Finite paths form a set, so \(D(a)\) exists.
It is the least predecessor-closed subset containing \(a\), and
\[
 D(a)=\{a\}\cup\bigcup_{bRa}D(b).
\]
No \(D(b)\) with \(bRa\) contains \(a\), since that would give a
finite \(R\)-cycle. The points on such a cycle form a nonempty
subset with no \(R\)-minimal member, contrary to well-foundedness.

A partial solution is a function on a predecessor-closed subset
of \(X\) satisfying the displayed recursion there.
Two partial solutions agree on their common domain.
Otherwise their set of disagreement has an \(R\)-minimal element.
All its predecessors have equal values in both solutions, so the
two sets of predecessor values are equal at that element,
contradicting disagreement.

Suppose a solution on \(D(b)\) exists for each \(bRa\).
It is unique by the preceding argument, so Replacement collects
these solutions without any choice of witnesses.
Their union is a function on \(\bigcup_{bRa}D(b)\).
Append the pair assigning to \(a\) the set of its predecessor values.
The appended point is new; the earlier domain is predecessor closed.
Thus no earlier recursion equation changes, and the extension is
a solution on \(D(a)\).
If any \(a\) had no such solution, an \(R\)-minimal member of
the set of failures would contradict this construction.
Therefore every \(a\) has a unique solution on \(D(a)\).
Replacement and union assemble them into the required function
\(\pi\) on \(X\); the agreement argument proves uniqueness.

Assume extensionality of \(R\).
Let \(B\) be the set of \(a\in X\) sharing their \(\pi\)-value
with a distinct point. If it is nonempty choose an \(R\)-minimal
\(a\in B\), and choose \(a'\neq a\) with \(\pi(a')=\pi(a)\).
For every \(bRa\), minimality gives \(b\notin B\).
Its image appears in \(\pi(a')\), so some predecessor of \(a'\)
has that image and must be \(b\).
Conversely, every predecessor of \(a'\) has the image of a
predecessor \(bRa\), which again determines that point uniquely.
Hence \(a,a'\) have the same predecessors, contradicting
extensionality. Thus \(\pi\) is injective.

Its range is transitive by the recursion equation.
The same equation and injectivity show
\(bRa\) exactly when \(\pi(b)\in\pi(a)\).
Any isomorphism onto a transitive set satisfies this equation:
each member of an image point is itself in the range and hence
is the image of a predecessor.
Uniqueness of the recursion therefore gives uniqueness of the
isomorphism and its transitive range.
\end{proof}
\end{annexitem}

\begin{annexitem}{annex:p06-finite-assignment-def}{thm:p06-finite-assignment-def}{Finite support of set satisfaction}
\begin{proof}
Fix nonempty \(A\), an instruction code \(c\), and \(m\geq1\)
containing all variable indices in that code.
For each instruction, form its truth set as a subset of \(A^m\).
Atomic instructions use coordinate equality or membership.
The false instruction gives the empty set.
Boolean instructions use the corresponding set operations.
An existential instruction at coordinate \(i\) includes a tuple
exactly when some update of that coordinate by an element of \(A\)
belongs to the earlier truth set; the universal instruction
uses every such update.
These are Separation definitions on the set \(A^m\).
Finite recursion along the instruction list gives a unique
evaluation table. Its construction and checking conditions are
uniform first-order formulas about set functions and codes.

Induction along the instruction list shows that two full
assignments agreeing below \(m\) give the same value, and
that this value is the one in the finite table.
For the quantifier steps, their updates at coordinate \(i<m\)
still agree on all coordinates used by the earlier instruction.
This verifies the induction at both quantifiers.

Every finite tuple extends to a full assignment: retain its
entries below \(m\) and use its zeroth entry on the remaining
natural numbers. Replacement forms this function.
Thus passing between full and finite assignments loses none of
the definable subsets used in \(\operatorname{Def}(A)\).
Codes and finite tuples form sets; Separation constructs their
uniform truth relation, using the finite evaluation formula.
For each code and tuple, Separation forms its output subset of
\(A\), and Replacement collects all these output subsets.
This is the complete definability family.
The empty base retains the separately specified value
\(\{\varnothing\}\).
All quantified evaluation structures have set domains.
\end{proof}
\end{annexitem}

\begin{annexitem}{annex:p06-hierarchy-finite-reflection}{thm:p06-hierarchy-finite-reflection}{Closing a level under the required witnesses}
\begin{proof}
Both hierarchies are increasing, transitive, and continuous at
limit ordinals. The constructible class is their union by
definition in the \(L\) case.
In the \(V\) case every set belongs to some level.
For completeness, the compatible-solution construction in
\cref{thm:p06-collapse-complete} also works with the value
\(\sup\{r(y)+1:y\in x\}\) on a transitive closure:
once predecessor values are given this is one uniquely specified
set value. Well-founded induction makes every value an ordinal,
and induction then gives \(x\in V_{r(x)+1}\).
This establishes the exhaustion property needed below.

Rewrite universal quantifiers using negation and existential
quantification, and close the fixed finite formula list under
subformulas. For each existential formula in this finite list
and each tuple of parameters from \(H_\alpha\), consider the
following bound. If the formula has a witness in \(H\), take
the least ordinal \(\xi\) such that a witness lies in \(H_\xi\);
otherwise use zero.
For each of the finitely many formulas this is a definable
functional rule on a set of parameter tuples.
Its definition uses that one formula's relativization to \(H\),
not a variable-code satisfaction predicate for the class.
Replacement bounds all these ordinals at once.

Starting above \(\beta\), define a strictly increasing sequence
\(\alpha_n\), choosing \(\alpha_{n+1}\) larger than the preceding
level and all its witness bounds.
Each choice is the uniquely specified supremum plus one,
so ordinary definable recursion on \(\omega\) suffices and
no Choice is used. Let \(\delta=\sup_n\alpha_n\).
It is a limit ordinal above \(\beta\).

Every finite tuple from \(H_\delta\) already lies in some
\(H_{\alpha_n}\), by continuity and finiteness.
If an existential formula from the list holds in \(H\) on
that tuple, the construction supplies a witness in
\(H_{\alpha_{n+1}}\subseteq H_\delta\).
Induction on formulas now proves the claimed equivalence.
Atomic formulas agree on their actual set arguments;
Boolean cases preserve equivalence.
For an existential formula, use the witness just obtained
for the forward direction and the induction hypothesis on
its matrix. The reverse direction uses the same matrix
hypothesis on any witness already in \(H_\delta\).
This covers every formula in the original list.
\end{proof}
\end{annexitem}

\begin{annexitem}{annex:p06-inner-l-zf}{thm:p06-inner-l-zf}{The ZF axioms one by one}
\begin{proof}
The class \(L\) is transitive and contains every ordinal by
\cref{thm:p06-constructible-class-foundation}.
In particular it contains the empty set and \(\omega\).
Every fixed finite list of its parameters lies in a common level.

Extensionality holds because members of constructible sets are
constructible. Two different constructible sets therefore have a
distinguishing member in the relativized domain.
The empty set is already present.
For Pairing, put \(x,y\) in a common \(L_\alpha\);
the subset of \(L_\alpha\) defined by \(z=x\lor z=y\)
belongs to \(L_{\alpha+1}\).
For Union, choose \(x\in L_\alpha\).
Transitivity ensures that every element of \(\bigcup x\)
lies in \(L_\alpha\). The formula
\(\exists y\in x\;(z\in y)\) defines that union over
\(L_\alpha\), so it belongs to \(L_{\alpha+1}\).
The ordinal \(\omega\), with the empty set and successor
operations just justified, witnesses Infinity.
For Foundation, a nonempty constructible set has an
ambient membership-minimal member. That member is constructible,
and its disjointness from the original set is a bounded
assertion, hence remains true in the transitive class.

For Separation, fix its standard formula \(\varphi\) and
constructible parameters \(a,\bar p\).
Use finite reflection to choose \(L_\delta\) containing
them and reflecting \(\varphi\) from \(L\).
Then
\[
 \{x\in a:\varphi^L(x,\bar p)\}
 =
 \{x\in L_\delta:x\in a\ \text{and}\
           (L_\delta,\in)\models\varphi(x,\bar p)\}.
\]
The right side is a definable subset of \(L_\delta\)
with its indicated parameters, and therefore belongs to
\(L_{\delta+1}\). This proves the given Separation instance.

For Replacement, suppose a fixed \(\varphi\) defines,
inside \(L\), a unique output for each member of
\(a\in L\), with parameters \(\bar p\in L\).
The ambient first-order formula
\(y\in L\land\varphi^L(x,y,\bar p)\) is functional on
the set \(a\).
Ambient Replacement produces its range \(b\).
At this point \(b\) is only known to be an ambient set
of constructible elements.
Assign to each member of \(b\) its least constructible
stage. Replacement bounds those stages, so \(b\) is
contained in some \(L_\eta\).
Choose a reflecting \(L_\delta\) above that bound and
containing \(a,\bar p\).
It defines \(b\) as the subset of its domain satisfying
\(\exists x\in a\,\varphi(x,y,\bar p)\).
Reflection rules out both missing and extra outputs.
Hence \(b\in L_{\delta+1}\), as required.

For the internal Power Set axiom, ambient Separation forms
\[
 p=\{y\in\mathcal P(a):y\in L\}.
\]
This is an ambient set.
Bounding the constructible stages of its members by Replacement,
choose \(L_\delta\) containing all those members and containing
\(a\). The bounded condition \(y\subseteq a\) is absolute
over this transitive level.
Consequently
\(p=\{y\in L_\delta:y\subseteq a\}\), a definable subset
of \(L_\delta\), and \(p\in L_{\delta+1}\).
These are exactly the subsets quantified over by the
relativized power-set statement.

All uses of Separation and Replacement above concern fixed
standard formulas. Each argument is a proof in \(\mathsf{ZF}\)
of the corresponding relativized axiom instance.
No single truth predicate for \(L\) was used.
\end{proof}
\end{annexitem}

\begin{annexitem}{annex:p06-l-self-construction}{thm:p06-l-self-construction}{Self-construction and the canonical well-order}
\begin{proof}
\emph{The hierarchy inside \(L\).}
The preceding theorem makes all set constructions of
\(\mathsf{ZF}\) available internally in \(L\).
It has the same ordinals and the same \(\omega\).
Finite tuples of constructible elements are constructible
by Pairing and Union, so finite assignments agree with
the ambient ones for a fixed set \(A\in L\).

Use the finite-assignment evaluation of
\cref{thm:p06-finite-assignment-def}.
Induction through a finite instruction code shows that its
truth over the set structure \((A,\in)\) is the same
inside \(L\) and outside it. Atomic membership and equality
are unchanged; every quantifier ranges over the same
elements of \(A\), all of which lie in \(L\).
Every finite-code definable subset is present internally by
Separation, and every internal definable subset has such
a code and finite parameters. Thus
\(\operatorname{Def}^L(A)=\operatorname{Def}(A)\).
The empty-base clause is also the same.
Transfinite induction, using the common successor
operation \(\operatorname{Def}\) and common unions at limits,
now identifies every internally computed \(L_\alpha\)
with its ambient value. Every member of \(L\) consequently
satisfies the internal assertion of constructibility.

\emph{A canonical order.}
Simultaneously with the levels construct well-orders \(W_\alpha\)
of \(L_\alpha\), each an initial-segment extension of the earlier ones.
Start with the empty order, and order
\(L_1=\{\varnothing\}\) uniquely. This handles the explicit
empty-base clause of \(\operatorname{Def}\).
For nonempty \(L_\alpha\), given \(W_\alpha\), list the
membership formulas by a fixed
natural-number coding, recording each formula's finite
parameter arity.
Order its definition codes first by the formula number,
then lexicographically by the parameter tuple in \(W_\alpha\).
For a fixed arity this finite lexicographic product is
well-ordered: minimize its first coordinate, then its next,
through its finitely many coordinates.
The ordered sum over the natural-number codes is therefore
a well-order of the full set of definition codes.

Each element of
\(L_{\alpha+1}\setminus L_\alpha\) has at least one such
definition code. Assign its least code.
Place these new elements after the old ones, in least-code order.
This specifies \(W_{\alpha+1}\) uniquely, without a choice
of definitions.
At a limit take the union of the preceding orders.
It is a linear order because each earlier carrier is an
initial segment. To see well-foundedness, take a nonempty
subset of the limit level and the least earlier stage
it meets. Its least element in that stage precedes
all its other members, including those at later stages.
Thus the union is again a well-order.

All code domains, least-code selections, and order unions
are sets at a set stage. The constructions use set-sized
satisfaction, not satisfaction for the universe.
Definable ordinal recursion therefore makes this a
parameter-free formula-defined class sequence.
The orders agree on overlaps, so they define one class
order of \(L\).

In \(\mathsf{ZF}+(V=L)\), any set lies in a constructible
level. Restrict that level's order to it by Separation.
Hence every set is well-orderable.
For a set of nonempty sets, well-order its union and assign
to each member its least element; Replacement forms a
choice function. This proves Choice from \(\mathsf{ZF}+(V=L)\).
Applying this derivation inside the already established
inner model \(L\) proves its Choice axiom.
Induction through the common levels compares the definition-code
orders, their least representatives, and the limit unions.
It gives \(W_\alpha^L=W_\alpha\) at every stage.
Thus the restriction of this canonical order to any constructible
set is itself a set in \(L\).
\end{proof}
\end{annexitem}

\begin{annexitem}{annex:p06-adequate-condensation}{thm:p06-adequate-condensation}{The finite-package collapse}
\begin{proof}
The membership relation on \(X\) is well-founded by ambient Foundation.
It is extensional: if two elements of \(X\) differ, transitivity
of \(L_\alpha\) supplies a distinguishing member in \(L_\alpha\);
elementarity supplies such a witness in \(X\).
Apply \cref{thm:p06-collapse-complete} and write \(M\) for the
transitive collapse. By elementarity and isomorphism,
\(M\models\Theta\).

We verify the absoluteness needed for this particular package.
The ordinal that \(M\) calls its least inductive set contains
every actual finite ordinal, by external induction.
If it were larger than the actual \(\omega\), transitivity
would put \(\omega\) itself in \(M\).
Pairing and Union compute the actual successors there, so
\(M\) would recognize that \(\omega\) is inductive,
contradicting leastness of its larger candidate.
Thus its natural numbers are exactly \(\omega\).
Pairing and Union also put every finite tuple of members
of \(M\) in \(M\). Formula codes and finite assignments over
a set \(A\in M\) therefore agree with the ambient ones.

The finite-assignment satisfaction clause of \(\Theta\)
has the same values in \(M\) and outside it.
Induct through the actual finite instruction code.
Atomic relations are absolute, and each quantifier ranges
over the same set \(A\), all of whose elements belong to \(M\).
The complete-output clause for \(\operatorname{Def}(A)\)
ensures both directions of the comparison:
every finite-code subset has an output in \(M\), and
every member of the internal definability family comes
from such a code and parameters.
Hence
\(\operatorname{Def}^M(A)=\operatorname{Def}(A)\),
including the empty-base clause.
This conclusion does not assume that \(M\) contains all
ambient subsets or all ambient infinite assignments.

Let \(\gamma=M\cap\operatorname{Ord}\).
It is an ordinal since \(M\) is transitive.
Pairing and Union give the actual successor of every
ordinal in \(M\), so \(\gamma\) is a limit ordinal.
For every \(\eta<\gamma\), the hierarchy-sequence clause
of \(\Theta\) supplies its internal levels through \(\eta\).
Their successor operations are the actual \(\operatorname{Def}\),
as just proved, and their limit unions are absolute.
External transfinite induction thus gives
\(L_\eta^M=L_\eta\).
The internal \(V=L\) assertion implies
\(M\subseteq\bigcup_{\eta<\gamma}L_\eta\).
Conversely \(L_\eta^M\) is an element of \(M\);
transitivity then gives \(L_\eta\subseteq M\).
Therefore \(M=L_\gamma\).

Bounded ordinalhood is absolute in both transitive sets.
The collapse consequently identifies \(X\cap\alpha\),
with its inherited membership order, with \(\gamma\).
The order type of a subset of \(\alpha\) is at most
\(\alpha\): in its increasing enumeration, the value
at position \(\xi\) is at least \(\xi\), by induction.
Hence \(\gamma\leq\alpha\).

Finally, \(L\) satisfies the construction clauses because
it satisfies \(\mathsf{ZF}\), and satisfies \(V=L\) by
\cref{thm:p06-l-self-construction}.
Apply the finite reflection schema to the fixed list
\(\Theta\), above any prescribed ordinal and above \(\omega\).
It gives a limit level satisfying the entire package.
This proves that the required stages are arbitrarily high.
\end{proof}
\end{annexitem}

\begin{annexitem}{annex:p06-gch-inner-l}{thm:p06-gch-inner-l}{The constructible cardinal bound}
\begin{proof}
Work first in \(\mathsf{ZF}+(V=L)\).
Choice is available by
\cref{thm:p06-l-self-construction}.

\emph{Sizes of levels.}
For every infinite set \(B\),
its definable subsets are indexed by countably many
formulas and their finite parameter tuples.
Thus
\[
 |\operatorname{Def}(B)|\leq
 \max(\omega,|B|).
\]
The same upper bound with \(\omega\) covers finite \(B\).
Transfinite induction gives
\[
 |L_\eta|\leq\max(\omega,|\eta|)
 \qquad\text{for every ordinal }\eta.
\]
The successor step uses the preceding definition-code bound.
At a limit, take the union of the earlier levels.
There are at most \(|\eta|\) of them and each has size at
most \(\max(\omega,|\eta|)\), so Choice and infinite-cardinal
multiplication give that same bound for their union.
These are the ordinary finite-tuple and indexed-union cardinal
bounds; no regularity of \(\eta\) is assumed.

\emph{Bounding the constructible stage of a subset.}
Fix an infinite cardinal \(\kappa\) and \(x\subseteq\kappa\).
Choose a sufficiently high limit \(\alpha\) satisfying
\(\Theta\), with \(x\in L_\alpha\) and
\(\kappa\subseteq L_\alpha\).
Such a level exists by
\cref{thm:p06-adequate-condensation}.

The structure \((L_\alpha,\in)\) is a set structure and
has a set-sized satisfaction relation.
For each existential formula, select the least witness
in the canonical well-order of \(L_\alpha\), when a
witness exists; use \(\varnothing\) otherwise.
This gives countably many finite-arity Skolem functions.
There is no use of a satisfaction predicate for the universe.

Start with \(X_0=\kappa\cup\{x\}\).
Given \(X_n\), adjoin the values of all these functions
on all finite tuples from \(X_n\), and call the result
\(X_{n+1}\).
Let \(X=\bigcup_{n<\omega}X_n\).
Each stage has cardinality at most \(\kappa\), since
\(\kappa\) is infinite and there are only countably
many functions with finite arities.
The union also has cardinality at most \(\kappa\).
Every finite tuple from \(X\) lies in one \(X_n\), so
every existential formula true on that tuple in
\(L_\alpha\) has a witness in \(X_{n+1}\).
Induction on formulas, with this witness condition at
the existential step, proves \(X\preccurlyeq L_\alpha\).

Its collapse is \(L_\gamma\) by the preceding lemma.
The collapse is a bijection, so
\(|L_\gamma|=|X|\leq\kappa\).
Since \(\gamma\subseteq L_\gamma\), it follows that
\(|\gamma|\leq\kappa\), and therefore \(\gamma<\kappa^+\).

The collapse fixes every \(\xi<\kappa\).
Indeed all members of such an ordinal are in \(X\),
and induction gives
\(\pi(\xi)=\{\pi(\zeta):\zeta<\xi\}=\xi\).
Every member of \(x\) lies below \(\kappa\), so
\(\pi(x)=x\) as well.
Consequently \(x\in L_\gamma\subseteq L_{\kappa^+}\).

\emph{The power-set cardinal.}
The argument applies to every subset of \(\kappa\);
under \(V=L\) there are no nonconstructible subsets to exclude.
Hence
\(\mathcal P(\kappa)\subseteq L_{\kappa^+}\).
The level bound gives
\(|\mathcal P(\kappa)|\leq\kappa^+\).
Cantor's diagonal argument gives
\(|\mathcal P(\kappa)|>\kappa\), and the definition of the
successor cardinal yields the reverse inequality
\(\kappa^+\leq|\mathcal P(\kappa)|\).
Therefore \(2^\kappa=\kappa^+\).
This proves GCH from \(\mathsf{ZF}+(V=L)\).

The earlier proofs establish inside \(L\) all axioms of
\(\mathsf{ZF}+(V=L)\) and Choice.
Relativizing the derivation just given therefore proves
GCH inside \(L\).
Together these give the claimed relativization of every
axiom of \(\mathsf{ZFC}+\mathsf{GCH}+(V=L)\).
Applying the finite proof-translation argument of
\cref{thm:ch-relative-consistency} yields its positive
relative-consistency implication without a cited
constructibility proof dependency.
\end{proof}
\end{annexitem}
